horizontal arrows - procedure steps

vertical arrows, for all the arrows, differntialte control related arrows, data flow related arrows

if thye correpond to differnt meaning, thye should use one type of arrow


data input/



Sectiion 5

Tables/figures shows most numbers, in your text, you dont have to enumaerate, you only need to do two things

1. idenify the point /conclusion or message you wan to pass across to the users e.g a data deltion methos is more useful on one paltform to another, you can point the user to the table

come up with a message, e,g overall

2. for this method, it is efecctive here thatn there

review the text tomake susre each number corresponds to a message i wan to deliver


bad example - all the chatbots fail to dlete user dats under some situations e,g chatboat a with this tuation, chatchbo c with thtis situatoo -  this is enumeration

dont start with enumerating numbers, start with a message, you dont have to enumerate all numbers


effectiveness of erasing data


it's not neccessary to list all these numbers

for the other olatofmrs, the results ae mostly cant't tell, with ranges of percent


need to make sure the number that should be htere, number, total, percent or ranges




dont make section 5 too long, bring in subheadings, if paragraph is too long, you can lost user, remember to highlight so it makes users understanding easy



CONSIDER THE REVIEWERS TO BE TOO GOOD AND NOTHNG

1.TOO GOOD, SOLVE IMPORTNAT PROBLRM AND DEMONSTRATE DOMAIN KNOWLEDGE
2. WTIE FOR LAYMANS TOO


WRITE FOR BOTH - HIGHLIGHT RESULT, MAKE THE PAPER EASIER TO AVIGATE

STRUCTURE MATTERS,
RECUE EFFORT FOR


HOW DO YOU DECIDE APPENDIX, REVIEW APPENDI, METHIDS AN PROCESS MATERS

DONT WORRY ABOUT APPNDINX YET, PUT ITIN THE RELVANT SECTION, WE CAN DECIDE WHICH ARE GOINF TO  APPENDINX LATER


CONVINCE REVIEWRS WITH LESS BACKGORUND, HIGHLITGHT IMPORTSANT THINGS

PIPELINE 
FINDING


WHAT IS MAJOR FINDING

ANYTHRING ELSE/FACORS HELPS EXPLAIN THE DIFFERNCE ACCROSS PLATFORM

PLATFOMR DESIGN/PRIVACY REGUAKLTION

THEY MAY NOT EXISTS, SO WE SHOILD TRY OUR BEST TON IDNEIFY THEM FOR AN EXTENSIVE DISCUSSION

----

Section 5 (the main message):
- Message first: start each paragraph with the point or conclusion, such as "a deletion method works on one platform and not on another". Point the reader to the table or figure for the numbers.
- No enumeration: don't list a number for every agent or method. "Chatbot A fails in this situation, chatbot C in that one" is the example he calls bad.
- Every number must serve a message: keep it only if a message needs it, and give it with its total, its percent, or a range. For the platforms where the result is mostly can't tell, give a range of percent and don't list them one by one.
- Length: keep Section 5 short, use subheadings, and highlight the key results.

Writing for reviewers:
- Write for expert reviewers and for laypeople too, and make the paper easy to navigate.
- Use less background and highlight the important things. Structure matters.
- Don't worry about the appendix yet. Put material in the relevant section, and we decide later what moves.

Findings and discussion:
- State the major finding clearly.
- Look for what explains the differences across platforms, such as platform design and privacy regulation. They may not exist, but we should try to identify them, because that gives the discussion its depth.

Figures: use horizontal arrows for procedure steps and vertical arrows for the rest. Give control arrows and data-flow arrows different looks, and use one arrow type per meaning.

This conflicts with what we just wrote. Our Section 5 is the enumeration he warns about in several places:
- 5.1, paragraph 2: the four method counts (6 of 13, 4 of 8, 2 of 7, 4 of 10).
- 5.1, paragraph 3: every agent's count.
- 5.2: "130 of 133 ChatGPT, 131 of 133 Claude, 57 of 133 Copilot", the per-code comparison, and the long documentation paragraph.

I haven't changed anything. A sensible order would be:
1. For each subsection, agree the one message it delivers.
2. Cut every number that doesn't support a message, and move the detail to the tables or provisional.md.
3. Add subheadings and highlights.


--



For the writing, I take it to mean:
- Surface reviewers need the message, the highlight and the navigation. Each paragraph opens with its point, defined terms are used the same way throughout, and subheadings let them find the result without reading every number.
- Deep reviewers need the proof that we know the domain and that the claim is solid. The numbers in tables and figures, the base for every count, and the links to the platform documentation and privacy rules serve them.

This matches the Section 5 note: the message in the text serves the first group, and the tables and figures with the full counts serve the second. It also explains why the notes want less background but more attention to what explains the platform differences.




he writing should follow a continuous methodological narrative, with paragraphs organized around distinct ideas and transitions connecting them. Headings should establish major conceptual divisions, not announce every procedure

let's answer the questions, what are the results available, whicj ones are importnat, how do we frame it, what do we want to communicate int his section and it's
  subsections, how do we wamt to drive the user, what value does it add to the paper or eesearch question, how will the reviewer see this, we need to answer all
  these questions to ge the best out pf this question


  ---

Results structures

- Yan (USENIX Sec '25, account deletion audit).
  - "Results and Findings" runs: Link Retrieval, Link Analysis, Method Classification, Verification.
  - The pilot study is organized by increasingly severe issues: link, then content, t
  - Its three RQs run availability, then compliance, then functionality.
- Hausladen (USENIX Sec '25, GPC).
  - Its Results open with "CCPA Opt Out Right Applicability: How Many Sites Sell or Share? Which Sites Are Subject?" before any compliance result.
  - It establishes the denominator first.
- Leave No Data Behind (PoPETs '22).
  - Results run by scope: individual data, account erasure, scope of erasure, access,
  - A separate "Further Insights" section holds the policies and the questionnaire.
- Khandelwal (USENIX Sec '24). RQ tags sit in the headings ("Data Practices in Privac
- Minaei, Kwesi and Malki. RQ tags sit in the headings, and Minaei adds a claim as each sub-heading.
- Setting the Course. It defines its properties first (completeness, correctness, conagainst them.

What this adds to my outline

1. Establish the denominator first. Hausladen's applicability step is the counterpartBefore saying how often deletion failed, say how many test cases could be judged at all (the fact reached the place we tested). That turns can't tell from an embarrassing 54.5% into the stated scope of the measurement.
2. A ladder from cheapest to hardest check. Yan's order is link, then content, then f the agent said, what the interface showed, what asking again returned, and what recall found.
3. RQ tags in the headings, as in Khandelwal, Minaei, Kwesi and Malki. This is the fo4 too, since your professor's headings stay and only the tag is added.
4. Keep secondary evidence separate. The documentation comparison goes in its own section or subsection, as in "Further Insights".

Revised outline
- Opening: the three outcomes, the lower-bound rule, and the RQ3/RQ4 map.
- 5.1 Overall Behavior of Agents [RQ3]:
  - (a) what could be judged;
  - (b) how often the fact came back among those;
  - (c) methods as a range;
  - (d) agents in a table.
- 5.2 Cross-Agent Difference [RQ3]: where each agent keeps the fact, the documentatio
- 5.3 What Users Can See: the ladder.
- 5.4 User Expectation vs. Behavior [RQ4]: the figure.

Decisions for you
1. Headline on all 865 or on the judgeable cases? On the judgeable cases it would read something like "of 394 test cases we could judge, the fact came back in 41%", which is larger and fairer to the
   measurement. On all 865 it stays at "at least 18.7%". I'd report both, with the jus provisional, since 61 test cases are placeholders.
2. RQ tags in the Section 5 and 4 headings: yes or no?
3. Merge "What Agents Say" into 5.3, or keep four subsections?
------

"Gemini permission cards" cant' tbe in the appendix, it is a finding too
--

) A one-paragraph recap of the findings from Sections 4
and 5.
2) Recommendations, one heading each, for example “State
what each deletion method deletes” and “Give users a
way to check”.
3) For platforms and regulators: deletion rights have to be
testable.
4) Limitations, short and last. The free-text reasons partici-
pants gave (not coded) go here as a note without counts

